\section*{Chapitre \ref{ch:iv1:logic-riddles} --- Énigmes de logique}
\markboth{\protect\omTocMark{}{Réponses~: Énigmes de logique}}{\protect\omTocMark{}{Réponses~: Énigmes de logique}}

\begin{solution}{iq:iv1:logic-riddles:1}
Deux pesées pour 8 boules, et 81 boules avec 4 pesées.

Une balance a trois issues\index{comptage d'information}~: plateau gauche plus lourd, plateau droit plus lourd, ou équilibre. Donc $k$ pesées distinguent au plus $3^k$ cas. On met 3 contre 3. Si un plateau est plus lourd, la boule est parmi ces 3. Sinon, elle est parmi les 2 restantes. Une pesée de plus la trouve alors parmi 3, en mettant 1 contre 1.

\iqwrong{Diviser par deux à chaque fois, comme si la balance n'avait que deux issues~: 4 contre 4, puis 2 contre 2, puis 1 contre 1. Cela demande 3 pesées, et n'utilise jamais l'issue d'équilibre.}
\end{solution}

\begin{solution}{iq:iv1:logic-riddles:2}
On allume la corde A par ses deux extrémités et la corde B par une seule. Quand A est entièrement consumée, au bout de 30 minutes, on allume l'autre extrémité de B. Elle se consume 15 minutes plus tard, soit 45 minutes au total.

Brûler une corde par les deux bouts divise le temps par deux, quelle que soit l'irrégularité de la corde\index{raisonnement par invariance}, car les deux flammes consument ensemble les 60 minutes de combustion. Quand A est consumée, la corde B a encore 30 minutes de combustion, et la seconde flamme divise ce temps par deux.

\iqwrong{Couper une corde en deux ou en quatre. La vitesse de combustion est irrégulière, donc la longueur ne dit rien du temps. Seuls les totaux de temps peuvent servir.}
\end{solution}

\begin{solution}{iq:iv1:logic-riddles:3}
On lance les deux sabliers. À 4 minutes, on retourne le sablier de 4 minutes. À 7 minutes, on retourne le sablier de 7 minutes. À 8 minutes, le sablier de 4 minutes est vide, et celui de 7 minutes a coulé depuis 1 minute~: on le retourne, et il coule alors pendant cette minute. On termine à 9 minutes.

L'astuce consiste à utiliser le sable déjà tombé\index{sablier}. Un sablier qui a coulé pendant $t$ minutes, une fois retourné, mesure exactement $t$ minutes. Il faut donc trouver un moment où un tel sablier donne la durée voulue~: ici, 1 minute après la marque des 8 minutes.

\iqwrong{Essayer de combiner seulement des multiples, comme $4a + 7b$, qui ne permet pas d'obtenir 9. On n'obtient 9 qu'en retournant un sablier à un moment où il est déjà en cours d'écoulement.}
\end{solution}

\begin{solution}{iq:iv1:logic-riddles:4}
On prend 1 pièce dans le sac 1, 2 dans le sac 2, et ainsi de suite jusqu'à 10 pièces dans le sac 10. Cela fait 55 pièces, qui pèseraient 550 grammes si toutes étaient authentiques. L'excès en grammes donne le numéro du sac lourd.

Une seule lecture d'une balance peut porter de nombreuses valeurs, donc on fait contribuer chaque sac d'une quantité différente\index{codage par les poids}. L'excès vaut $k$ lorsque le sac $k$ est le faux, et cette valeur est différente pour chacun des dix cas.

\iqwrong{Peser un sac à la fois, ce qui demande jusqu'à neuf utilisations de la balance. Une balance donne un nombre, et non une réponse oui ou non.}
\end{solution}

\begin{solution}{iq:iv1:logic-riddles:5}
On prend un A de plus dans son flacon, de sorte que l'on ait 2 A et 2 B. On coupe chacune des quatre pilules en deux, et l'on met une moitié de chacune dans le tas du jour, et l'autre moitié dans celui du lendemain.

On ne peut pas distinguer les pilules, donc on ne peut pas sélectionner selon le type, mais on peut rendre la sélection identique pour chaque pilule\index{symétrisation}. Les quatre moitiés du jour représentent la moitié de 2 A et de 2 B, soit exactement un A et un B. Les quatre moitiés du lendemain forment la même dose. Aucune pilule n'est perdue.

\iqwrong{Jeter les pilules et recommencer, ce qui ne répond pas à la question, ou en prendre deux parmi les trois, ce qui risque une mauvaise dose.}
\end{solution}

\begin{solution}{iq:iv1:logic-riddles:6}
Une seule coupure suffit, au troisième maillon (ou, par symétrie, au cinquième), ce qui laisse des morceaux de 1, 2 et 4 maillons. Avec 2 coupures, une chaîne de 23 maillons fonctionne.

Au jour $d$, l'ouvrier doit avoir exactement $d$ maillons, donc les morceaux doivent permettre d'obtenir chaque total de 1 à 7, en utilisant chaque morceau au plus une fois. Des morceaux de 1, 2 et 4 le permettent, comme en binaire\index{représentation binaire}. Ainsi, le deuxième jour, on reprend le 1 et l'on remet le 2.

Avec 2 coupures, on obtient deux maillons isolés, qui permettent d'atteindre 2. Le morceau suivant peut alors être de 3, ce qui donne 5, puis 6, ce qui donne 11, et enfin 12, ce qui donne 23. Donc $1 + 1 + 3 + 6 + 12 = 23$. Essayer toutes les paires de coupures confirme qu'aucune chaîne plus longue ne fonctionne.

\iqwrong{Trois coupures, aux maillons 2, 4 et 6, qui laissent sept maillons isolés à remettre un par jour. Cela ignore la monnaie rendue, qui est le cœur de l'énigme.}
\end{solution}

\enlargethispage{-3\baselineskip}%
\begin{solution}{iq:iv1:logic-riddles:7}
Non pour les coins opposés. Oui pour deux cases quelconques de couleurs opposées.

On colorie l'échiquier de la manière habituelle. Un domino recouvre une case noire et une case blanche\index{argument de coloriage}. Les coins opposés ont la même couleur, donc le reste comporte 30 cases d'une couleur et 32 de l'autre, et les dominos ne peuvent pas le recouvrir.

Pour deux cases de couleurs opposées, il existe une belle preuve. On trace un chemin fermé passant par les 64 cases, en allant de case voisine en case voisine, par exemple une forme de peigne. Retirer deux cases de couleurs opposées coupe le chemin en deux morceaux, chacun comportant un nombre pair de cases, car les couleurs alternent le long du chemin. Chaque morceau est recouvert par des dominos posés le long de lui. Une recherche exhaustive sur un échiquier 4 par 4 le confirme.

\iqwrong{Essayer de placer les dominos un à un, rester bloqué, et conclure que cela semble impossible. Cela ne prouve rien. Un décompte des couleurs le démontre en une ligne.}
\end{solution}

\begin{solution}{iq:iv1:logic-riddles:8}
L'un annonce toujours la couleur qu'il voit sur l'autre. L'autre annonce toujours le contraire de ce qu'il voit. Exactement l'un des deux a raison.

Si les chapeaux sont identiques, le premier a raison. S'ils diffèrent, c'est le second qui a raison\index{enigme@énigme des chapeaux}. Ce sont les deux seuls cas, donc l'un d'eux gagne toujours. Chacun a raison une fois sur deux, comme il se doit, puisqu'il ne sait rien de son propre chapeau. La stratégie ne change pas cela. Elle fait en sorte que les bonnes réponses ne coïncident jamais.

\iqwrong{Affirmer que c'est impossible, puisque chaque réponse est un pile ou face. C'est vrai pour chaque réponse prise isolément. Une stratégie peut corréler les réponses.}
\end{solution}

\begin{solution}{iq:iv1:logic-riddles:9}
On allume l'interrupteur 1 pendant dix minutes, on l'éteint, puis on allume l'interrupteur 2 et l'on entre dans la pièce. L'ampoule allumée correspond à l'interrupteur 2, l'ampoule éteinte et chaude à l'interrupteur 1, et l'ampoule éteinte et froide à l'interrupteur 3.

Un simple regard sur un état allumé ou éteint ne donne que deux valeurs, ce qui ne suffit pas pour trois interrupteurs\index{comptage d'information}. Une ampoule à filament possède un état supplémentaire, sa température. Il faut un troisième signal, et la chaleur en est un que l'on peut produire à l'avance.

\iqwrong{Affirmer que c'est impossible, car il n'y a qu'un seul passage et deux états seulement. Se demander ce qui peut encore différer d'une ampoule. Un recruteur peut vouloir voir le modèle remis en cause, et pas seulement le calcul.}
\end{solution}

\begin{solution}{iq:iv1:logic-riddles:10}
On prend un fruit dans la boîte étiquetée «~mixte~». Elle ne contient en réalité qu'un seul type de fruit, celui que l'on tire. Si l'on tire une pomme, cette boîte contient des pommes~; la boîte étiquetée «~oranges~» ne peut donc contenir ni des oranges ni des pommes, elle est donc mixte, et la boîte étiquetée «~pommes~» contient des oranges.

Toutes les étiquettes sont fausses, donc seules deux dispositions sont possibles\index{dérangements}. La boîte étiquetée «~mixte~» ne contient que des fruits d'une seule espèce dans les deux cas, donc un seul tirage dans celle-ci suffit à identifier la disposition. Un tirage dans la boîte étiquetée «~pommes~» peut donner une orange dans les deux dispositions, et dans ce cas il n'apprend rien.

\iqwrong{Choisir une boîte au hasard. Un fruit tiré d'une boîte qui peut être mélangée ne dit pas ce qu'elle contient.}
\end{solution}

\begin{solution}{iq:iv1:logic-riddles:11}
Le temps minimal est de 17 minutes. Les personnes 1 et 2 traversent (2 minutes), 1 revient (1), les personnes 5 et 10 traversent (10), 2 revient (2), puis 1 et 2 traversent (2).

Le plan naturel consiste à faire escorter tout le monde par la personne la plus rapide, ce qui donne $10 + 1 + 5 + 1 + 2 = 19$. Le meilleur plan fait traverser ensemble les deux plus lentes\index{argument d'échange}, de sorte que le 5 ne coûte rien de plus, au prix de devoir faire revenir la personne 2 une fois. Une recherche sur tous les états donne un minimum de 17.

\iqwrong{Donner 19, en supposant que la personne la plus rapide doit toujours porter la lampe. Regrouper les deux plus lentes fait gagner davantage que l'escorte rapide.}
\end{solution}

\begin{solution}{iq:iv1:logic-riddles:12}
On remplit le bidon de 5, on verse dans celui de 3, il reste 2. On vide le 3 et l'on y verse les 2. On remplit le 5, et l'on complète le 3 avec 1 litre, ce qui laisse 4. Avec des bidons de 6 et 10, on ne peut pas obtenir 3.

Toute quantité obtenable est une combinaison des deux capacités, donc un multiple de leur plus grand commun diviseur\index{identité de Bézout}. Et l'on peut atteindre chaque multiple de ce diviseur jusqu'à la plus grande capacité. Pour 3 et 5, le diviseur vaut 1. Pour 6 et 10, il vaut 2, et seules les quantités paires apparaissent donc. Une recherche sur tous les états des bidons confirme les quantités $\{0, 2, 4, 6, 8, 10\}$.

\iqwrong{Essayer des coups au hasard jusqu'à obtenir la réponse. Énoncer la règle, qui indique aussi quand s'arrêter.}
\end{solution}

\begin{solution}{iq:iv1:logic-riddles:13}
Non. La parité du nombre de cases noires ne change jamais, et elle vaut 1 au départ.

Un coup retourne 8 cases\index{invariant de parité}. Si $k$ d'entre elles étaient noires, le nombre de cases noires change de $8 - 2k$, qui est pair. Le compte reste donc impair, et il ne peut pas valoir 0. Une longue suite de coups aléatoires sur une grille confirme que le compte reste impair à chaque étape.

\iqwrong{Essayer de corriger l'unique case noire en retournant sa ligne. Cela rend la ligne noire sauf une case, et le compte passe à 7~: de nouveau impair.}
\iqfollow{Et sur une grille 7 par 7~?}{Un coup retourne maintenant 7 cases, donc le compte ne sert plus à rien. On utilise plutôt un bloc 2 par 2 quelconque~: un coup en retourne 0 ou 2 de ses cases, donc sa parité ne change jamais, et un bloc qui contient la case noire reste impair. La réponse reste donc non, quelle que soit la taille de la grille.}
\end{solution}

\begin{solution}{iq:iv1:logic-riddles:14}
Quatre-vingt-dix-neuf peuvent être sûrs, et la dernière personne du fond a une chance sur deux. Celle du fond dit «~noir~» si le nombre de chapeaux noirs devant elle est impair, et «~blanc~» sinon.

Le plan encode un bit, la parité des 99 chapeaux\index{codage par parité}. La personne suivante voit 98 de ces chapeaux. En comparant leur parité avec celle qui a été annoncée, elle connaît avec certitude sa propre couleur. Chacun à son tour connaît la parité annoncée, a entendu toutes les réponses venues de derrière, et voit le reste~; il peut donc déduire la couleur de son propre chapeau. Une simulation de 300 lignes aléatoires ne trouve aucune erreur parmi ces 99.

Donc, en moyenne, $99.5$ réponses sont justes. On ne peut pas faire mieux que la chance sur deux de cette personne, puisque personne n'a vu son chapeau.

\iqwrong{Associer les personnes deux à deux, pour que chacune annonce la couleur de celle de l'autre, et sauver ainsi environ la moitié à coup sûr. Cela gâche le fait que chacun entend tout, et qu'un seul bit suffit pour tous.}
\iqfollow{Et s'il y a trois couleurs~?}{On utilise la somme des numéros de couleur modulo 3, et on l'annonce sous la forme de l'un de trois mots. Toujours 99 sont sûrs. En général, une annonce porte une somme modulo le nombre de couleurs.}
\end{solution}

\begin{solution}{iq:iv1:logic-riddles:15}
Il faut 10 goûteurs pour un tour, et 7 goûteurs pour deux tours.

Un tour~: on écrit le numéro de chaque bouteille en binaire\index{codage binaire}, sur 10 chiffres, car $2^{10} = 1024 \ge 1000$. Le goûteur $k$ boit dans chaque bouteille dont le $k$-ième chiffre vaut 1. L'ensemble des goûteurs touchés donne les chiffres de la bouteille empoisonnée. On ne peut pas faire avec 9 goûteurs, car 9 goûteurs ne produisent que 512 résultats distincts.

Deux tours~: chaque goûteur a maintenant trois résultats possibles~: touché au premier tour, touché au second tour, ou jamais. On écrit donc les bouteilles en base 3, sur 7 chiffres, puisque $3^6 = 729 < 1000 \le 3^7 = 2187$. Un chiffre 0 signifie que le goûteur boit au premier tour, un chiffre 1 qu'il boit au second tour, et un chiffre 2 qu'il ne boit jamais. Un goûteur touché au premier tour n'a pas besoin de reboire au second tour, car ce chiffre est déjà connu.

En général, avec $r$ tours, $t$ goûteurs traitent $(r + 1)^t$ bouteilles.

\iqwrong{Utiliser 1000 goûteurs, ou une recherche dichotomique, qui prend 10 jours. Le test de groupe en parallèle consiste à concevoir un seul test dont les résultats encodent plusieurs bits à la fois.}
\iqwhy{C'est l'idée derrière le test de groupe et l'isolation rapide de pannes~: on met beaucoup de choses dans chaque test, et le motif de résultats de chacune sert d'adresse.}
\end{solution}

\begin{solution}{iq:iv1:logic-riddles:16}
Prendre n'importe quelles 10 pièces comme premier groupe, et les retourner toutes. Les deux groupes ont alors le même nombre de faces.

Supposons que les 10 choisies comptent $h$ faces\index{raisonnement par invariant}. Les 90 autres en comptent donc $10 - h$. Après avoir retourné vos 10 pièces, elles en comptent $10 - h$ également. Peu importe le nombre de faces que vous avez choisies, et c'est pourquoi vous n'avez pas besoin de les voir. Une simulation avec des dispositions aléatoires donne l'égalité à chaque fois.

\iqwrong{Dire qu'il faut toucher les pièces, ou chercher une façon de les répartir également. Le nombre nécessaire est la taille du groupe, et il découle du total donné, pas de la position des pièces.}
\end{solution}

\begin{solution}{iq:iv1:logic-riddles:17}
Non. Les deux coins ont la même couleur, mais un parcours du cavalier sur les 64 cases comporte 63 coups, et il doit finir sur la couleur opposée à celle du départ.

Chaque coup du cavalier change la couleur de la case\index{argument de coloriage}. Après 63 coups, nombre impair, on se retrouve sur la couleur opposée. Deux coins opposés de l'échiquier ont la même couleur, car $(0 + 0)$ et $(7 + 7)$ sont tous deux pairs. C'est donc impossible.

L'argument interdit seulement la même couleur, si bien qu'une arrivée sur l'autre couleur n'est pas exclue par cet argument. Il ne prouve pas à lui seul qu'un tel parcours existe. Sur de petits échiquiers, on peut tester la règle par une recherche exhaustive~: sur un échiquier 4 par 4, aucun parcours ne relie deux coins opposés, alors que sur un échiquier 5 par 5, avec 24 coups, un tel parcours existe.

\iqwrong{Essayer de dessiner un parcours et échouer. Un échec à trouver un parcours ne prouve rien, et sur de plus grands échiquiers on n'en viendrait pas à bout. La parité tient en une ligne.}
\end{solution}

\begin{solution}{iq:iv1:logic-riddles:18}
On raisonne par récurrence sur $n$. Un échiquier 8 par 8 utilise 21 pièces.

Pour $n = 1$, un échiquier 2 par 2 privé d'une case est une pièce en L. Prenons maintenant un échiquier $2^n$ par $2^n$ privé d'une case. On le découpe en quatre quadrants\index{démonstration par récurrence}. La case manquante se trouve dans l'un des quadrants. On place une pièce en L au centre de l'échiquier, de sorte qu'elle prenne une case dans chacun des trois autres quadrants. Chacun des quatre quadrants a alors exactement une case manquante, et chacun est un échiquier plus petit du même type.

Le compte est $(64 - 1)/3 = 21$, et $4^n - 1$ est toujours divisible par 3. Un programme réalisant cette construction, pour chacune des 64 positions du trou, recouvre tout avec 21 pièces de trois cases.

\iqwrong{Supposer que le seul comptage donne la réponse. Il montre seulement que les nombres s'accordent, et un coloriage pourrait encore interdire certaines positions. La récurrence donne un vrai recouvrement.}
\end{solution}

\begin{solution}{iq:iv1:logic-riddles:19}
Pas avec 13, 15 et 17. Avec 13, 16 et 16, oui.

Quand deux se rencontrent, un des comptes augmente de 2 et deux diminuent de 1\index{invariant modulo 3}. La différence entre deux comptes change donc de 0 ou de $\pm 3$, et elle reste inchangée modulo 3. Si tous étaient d'une même couleur, deux comptes vaudraient 0, avec une différence de 0. Dans 13, 15, 17, les restes sont 1, 0, 2, tous différents, si bien que deux comptes ne peuvent jamais être égaux modulo 3.

Avec 13, 16 et 16, deux comptes sont égaux. Faisons-les se rencontrer~: chaque rencontre entre vert et bleu fait diminuer les deux, et 16 rencontres transforment tous les autres en rouge. Une recherche sur tous les états atteignables confirme les deux réponses.

\iqwrong{Penser à la parité, comme pour tant d'invariants. Une rencontre change les comptes de 1 et de 2, donc la parité n'est pas conservée, et le bon module est 3.}
\end{solution}

\begin{solution}{iq:iv1:logic-riddles:20}
Ce ne peut pas être 0, car la parité de la somme ne change jamais et elle part impaire. Le second jeu se termine toujours sur $11! - 1 = 39\,916\,799$, quel que soit l'ordre.

D'abord, $|a - b|$ a la même parité que $a + b$\index{invariant de parité}, donc chaque coup conserve la parité du total. Le total part de 55, qui est impair, donc le dernier nombre est impair.

Ensuite, remarquons que $a + b + ab = (a + 1)(b + 1) - 1$\index{invariant produit}. Le produit des $x + 1$ sur tous les nombres du tableau ne change donc pas. Il vaut au départ $2 \cdot 3 \cdots 11 = 11!$, et à la fin il vaut le dernier nombre plus un. Des ordres aléatoires dans une simulation confirment ces deux résultats.

\iqwrong{Essayer quelques ordres pour chercher une régularité. Tout invariant se trouve plus facilement en se demandant ce que préserve un coup.}
\end{solution}

\begin{solution}{iq:iv1:logic-riddles:21}
Il faut 148 comparaisons pour 100 nombres. En général, c'est $\lceil 3n/2 \rceil - 2$.

On compare les nombres deux à deux\index{comptage des comparaisons}. Chaque paire utilise une comparaison, puis le plus grand de chaque paire dispute le maximum, et le plus petit de chaque paire dispute le minimum. Pour $n = 100$~: 50 comparaisons par paires, puis 49 pour le maximum, et 49 pour le minimum, soit 148. Un comptage sur des tableaux réels de tailles variées confirme la formule.

Pourquoi il est difficile de faire mieux~: pour conclure, $n - 1$ nombres doivent avoir perdu une fois, et $n - 1$ doivent avoir gagné une fois, donc il faut $2n - 2$ unités d'information. Une comparaison entre deux nombres encore intacts en apporte 2, et toute comparaison ultérieure en apporte au plus 1 face à un adversaire. Il y a au plus $n/2$ comparaisons du premier type.

\iqwrong{Utiliser $198$ comparaisons, ou affirmer qu'il est bon de trier d'abord. Trier coûte environ $n\log_2 n$. Un seul parcours pour les deux extrêmes ignore ce que chaque comparaison peut apprendre.}
\end{solution}

\enlargethispage{-3\baselineskip}%
\begin{solution}{iq:iv1:logic-riddles:22}
Les âges sont 2, 2 et 9, et le numéro de la maison est 13.

On liste les triplets de produit 36. Il y en a huit~: $(1, 1, 36)$, $(1, 2, 18)$, $(1, 3, 12)$, $(1, 4, 9)$, $(1, 6, 6)$, $(2, 2, 9)$, $(2, 3, 6)$, $(3, 3, 4)$. Leurs sommes valent 38, 21, 16, 14, 13, 13, 11 et 10. Le recenseur voit le numéro, et ne peut toujours pas trancher, donc il faut une somme qui apparaît deux fois\index{connaissance comme information}~: 13. Alors «~le plus âgé~» écarte $(1, 6, 6)$, où les deux plus âgés ont le même âge.

\iqwrong{S'arrêter lorsque la somme vous semble ambiguë. L'information tient au fait que le recenseur n'a pas pu décider, ce qui est une affirmation sur son point de vue, non sur le vôtre.}
\end{solution}

\begin{solution}{iq:iv1:logic-riddles:23}
Le chapeau de devant est rouge.

Il n'y a que 2 chapeaux blancs, donc si la personne du fond en voyait deux blancs, elle saurait que le sien est rouge\index{connaissance comme information}. Elle dit ne pas savoir, donc le milieu et le devant ne sont pas tous deux blancs. La personne du milieu le sait. Si le chapeau de devant était blanc, la personne du milieu le verrait et saurait que le sien doit être rouge. Or elle ne sait pas, donc le chapeau de devant n'est pas blanc.

La personne de devant ne voit personne, et raisonne à partir de ce que les deux autres ont dit. Chaque «~je ne sais pas~» est une information sur ce qu'ils ont vu.

\iqwrong{Dire que la personne de devant n'a aucune information parce qu'elle ne voit rien. Chaque affirmation réduit les mondes possibles, de sept à ceux où le chapeau de devant est rouge.}
\end{solution}

\begin{solution}{iq:iv1:logic-riddles:24}
On prend $2^{k - 1}$ pièces de la pile $k$~: 1, 2, 4, jusqu'à 512, soit 1023 pièces. L'excédent en grammes, écrit en binaire, donne la liste des piles truquées.

Chaque sous-ensemble de piles doit donner un excédent différent\index{codage binaire}. Les puissances de 2 le font, car chaque nombre a un développement binaire unique. Les 1024 sous-ensembles possibles demandent 1024 excédents distincts, à partir de 0, donc un total d'au moins 1023 grammes supplémentaires doit être possible, ce qui impose au moins 1023 pièces, et les puissances de 2 atteignent cette borne. Une vérification exhaustive des 1024 sous-ensembles confirme.

\iqwrong{Utiliser 1, 2, 3, jusqu'à 10, comme pour une seule pile truquée. Alors 3 pourrait désigner la pile 3, ou les piles 1 et 2 ensemble.}
\end{solution}

\begin{solution}{iq:iv1:logic-riddles:25}
Demandez à un garde~: «~Répondriez-vous oui si je vous demandais si le chemin de gauche est sûr~?~» Un oui signifie que la gauche est sûre, dans tous les cas.

Vérifions. Si la gauche est sûre et que le garde dit la vérité, la réponse directe serait oui, et ils disent oui. Si la gauche est sûre et que le garde ment, sa réponse directe serait non, et il ment là-dessus, donc il dit oui. Si la gauche n'est pas sûre, le garde qui dit la vérité répond non, et la réponse directe du menteur serait oui, sur quoi il ment, donc il dit non\index{double négation}.

Une autre forme interroge l'autre garde~: «~L'autre garde affirmerait-il que la gauche est sûre~?~» Ici la réponse est toujours fausse, donc un non signifie que la gauche est sûre. Un examen exhaustif des quatre cas confirme les deux formes.

\iqwrong{Poser une question qui identifie le menteur, par exemple «~Deux plus deux font-ils quatre~?~» Cela consomme l'unique question dont on dispose pour autre chose que la route.}
\end{solution}

\begin{solution}{iq:iv1:logic-riddles:26}
L'énoncé 99 est le seul vrai. Avec «~au moins $k$~», les énoncés 1 à 50 sont vrais, et les autres sont faux.

Au plus un énoncé «~exactement~» peut être vrai, car ils se contredisent deux à deux. Si aucun n'était vrai, les 100 seraient faux, ce qui rendrait l'énoncé 100 vrai~: une contradiction. Donc exactement un est vrai, et 99 sont faux, ce qui fait de l'énoncé 99 celui qui est vrai\index{auto-référence}.

Pour «~au moins $k$~», supposons que $f$ énoncés soient faux. Les énoncés vrais sont exactement ceux pour lesquels $k \le f$, il y en a donc $f$. D'où $100 - f = f$, et $f = 50$. Une recherche sur tous les nombres possibles d'énoncés faux montre que chaque version a une seule solution.

\iqwrong{Choisir l'énoncé 100, parce que c'est le dernier. S'il était vrai, les 100 seraient faux, y compris lui-même.}
\end{solution}

\begin{solution}{iq:iv1:logic-riddles:27}
Utiliser un schéma fixe~: peser $\{1, 2, 3, 10\}$ contre $\{4, 5, 6, 11\}$, puis $\{1, 2, 3, 11\}$ contre $\{7, 8, 9, 10\}$, et enfin $\{1, 4, 7, 10\}$ contre $\{2, 5, 8, 12\}$. Treize billes ne peuvent pas être traitées ainsi.

Chaque bille a un profil sur les trois pesées, $+$, $-$ ou $0$ pour le plateau gauche, le plateau droit ou hors balance. Si la bille est lourde, ce profil est le résultat. Si elle est légère, c'est l'opposé de ce profil\index{codage ternaire}. Donc les 24 cas doivent donner des résultats différents, et le profil de chaque bille ne doit pas être l'opposé de celui d'une autre. Dans ce schéma, la bille 1 est $(+, +, +)$, la bille 2 est $(+, +, -)$, la bille 10 est $(+, -, +)$ et la bille 11 est $(-, +, 0)$. Un examen exhaustif confirme que les 24 résultats sont tous différents et qu'aucun n'est équilibré sur les trois pesées.

Voyons maintenant pourquoi 13 échoue~: $27 - 1 = 26$ résultats et $13 \cdot 2 = 26$ cas semblent s'ajuster exactement. Mais regardons la première pesée. Plaçons $a$ billes sur chaque plateau. Si la balance penche, il reste $2a$ cas, et deux pesées ont 9 résultats, donc $a \le 4$. Si elle est équilibrée, il reste $13 - 2a$ billes, donc $2(13 - 2a)$ cas, qui doit être au plus 9, ce qui donne $a \ge 5$. Ces deux conditions se contredisent.

\iqwrong{Croire que $3^3 = 27$ règle le compte. Ce compte est une borne. Il faut aussi une première étape qui répartit les cas de façon égale, et l'exigence d'un nombre égal de billes sur les deux plateaux l'exclut pour 13.}
\end{solution}

\begin{solution}{iq:iv1:logic-riddles:28}
Non pour 10 sur 10, et oui pour 4 sur 10. Un échiquier de $m$ sur $n$ peut être couvert exactement lorsque 4 divise $m$ ou $n$.

Colorions la case en $(x, y)$ avec $(x + y) \bmod 4$\index{argument de coloriage}. Chaque brique droite couvre une case de chacune des quatre couleurs, quelle que soit son orientation. Sur un échiquier de 10 sur 10, les effectifs des couleurs sont 25, 26, 25 et 24, qui ne sont pas égaux. Il n'existe donc aucun pavage.

Pour la règle générale, utilisons un coloriage plus fort. Plaçons le nombre $i^{x + y}$ sur chaque case, avec $i^2 = -1$. Une brique porte quatre puissances consécutives de $i$, dont la somme est 0\index{racines de l'unité}. La somme sur un échiquier couvert est donc 0. Or la somme sur l'échiquier vaut
\[
\Bigl(\sum_{x = 0}^{m - 1} i^x\Bigr)\Bigl(\sum_{y = 0}^{n - 1} i^y\Bigr),
\]
et une somme de puissances consécutives de $i$ n'est nulle que lorsque sa longueur est un multiple de 4. Donc 4 doit diviser $m$ ou $n$. Et si c'est le cas, disposer les briques en rangées suffit. Une recherche exhaustive pour toutes les tailles jusqu'à 7 sur 7 confirme cette règle.

\iqwrong{Utiliser le nombre total de cases. Un échiquier de 10 sur 10 a 100 cases, qui est divisible par 4, et il ne peut tout de même pas être couvert. Un décompte n'est qu'une première vérification.}
\end{solution}

\begin{solution}{iq:iv1:logic-riddles:29}
Pour 64 nombres, 68 comparaisons, contre $63 + 62 = 125$ pour deux parcours. En général, c'est $n + \lceil \log_2 n \rceil - 2$.

Organisons un tournoi à élimination directe\index{méthode du tournoi}~: $n - 1$ comparaisons, qui trouvent le plus grand. Le deuxième plus grand a nécessairement perdu directement contre le plus grand, car il ne pouvait perdre contre personne d'autre. Il fait donc partie des nombres que le vainqueur a battus, dont il y a $\lceil \log_2 n \rceil$. Cherchons le plus grand parmi eux avec $\lceil \log_2 n \rceil - 1$ comparaisons supplémentaires. Pour $n = 64$, cela fait $63 + 5$. Une exécution sur des tableaux aléatoires donne le deuxième plus grand correctement, avec, dans le pire cas, exactement 68 comparaisons.

Pourquoi cela doit-il être ainsi~: il faut savoir que chaque nombre, sauf les deux premiers, n'est pas le deuxième. La seule preuve qu'un nombre n'est pas le deuxième est qu'il a perdu contre autre chose que le plus grand. Sinon, il pourrait encore être le deuxième. Donc $n - 2$ nombres ont besoin d'une défaite contre un nombre qui n'est pas le plus grand. En plus, le plus grand doit gagner ses propres comparaisons, et un adversaire peut faire en sorte qu'il en gagne au moins $\lceil \log_2 n \rceil$, dont aucune ne compte pour le premier groupe\index{argument de l'adversaire}.

\iqwrong{Trier, ou retirer le plus grand puis reparcourir la liste, pour 125 comparaisons. Ce second parcours répète des comparaisons dont on connaît déjà le résultat.}
\end{solution}

\begin{solution}{iq:iv1:logic-riddles:30}
Numérotons les personnes de 0 à $n - 1$, et les couleurs de 0 à $n - 1$. La personne $k$ annonce la couleur qui ferait que la somme de toutes les couleurs soit égale à $k$ modulo $n$. Exactement une est juste.

Soit $S$ la vraie somme, modulo $n$\index{énigmes de chapeaux@énigmes de chapeaux}\index{codage modulaire}. La personne $k$ voit tous les chapeaux sauf le sien, et suppose que le total vaut $k$. La personne $k$ a donc raison exactement lorsque $S = k$. Comme $S$ prend une seule valeur de $0$ à $n - 1$, exactement une personne a raison.

Chaque personne a raison avec une probabilité de $1/n$, comme il se doit, car chaque chapeau est indépendant de ce qu'elle voit. Le nombre total de bonnes réponses du groupe a toujours une moyenne de 1. Sans plan, ce nombre varie, et il peut valoir 0. Le plan supprime toute la variance. Un examen exhaustif de toutes les dispositions pour $n$ jusqu'à 5 trouve exactement une bonne réponse à chaque fois.

\iqwrong{Croire qu'un plan doit augmenter la chance de chacun. Il ne le peut pas, il doit seulement modifier la façon dont les bonnes réponses se recoupent.}
\iqwhy{La diversification fonctionne de la même façon. On ne peut pas modifier le gain espéré de chaque pari, mais on peut organiser la corrélation pour qu'exactement un paye dans chaque état.}
\end{solution}

\begin{solution}{iq:iv1:logic-riddles:31}
Ils peuvent gagner avec une probabilité de $3/4$. Le plan~: si l'on voit deux chapeaux de la même couleur, on annonce la couleur opposée. Si l'on voit deux chapeaux de couleurs différentes, on passe.

Il y a huit issues. Dans les 6 où deux chapeaux ont une couleur et le troisième l'autre, la personne au chapeau isolé voit deux chapeaux égaux et annonce correctement, et les deux autres voient des chapeaux différents et passent. C'est une victoire\index{code de Hamming}. Dans les deux où les trois chapeaux sont identiques, les trois annoncent faux, ce qui est une défaite.

Chaque annonce reste juste une fois sur deux. Sur les huit issues, le plan fait 12 annonces~: 6 sont justes, une dans chaque issue gagnante, et 6 sont fausses, toutes dans les deux issues perdantes. Le plan concentre les mauvaises réponses dans quelques issues.

Aucun plan ne fait mieux. Les bonnes et les mauvaises annonces sont en nombre égal pour tout plan. Chaque issue gagnante exige au moins une bonne annonce, et chaque issue perdante compte au plus 3 mauvaises annonces. Donc victoires $\le 3 \times$ défaites, ce qui donne au plus $6$ sur 8. Une recherche sur tous les $81^3$ plans confirme un maximum de 6.

\iqwrong{Dire que la probabilité est $1/2$, parce que toute annonce a une probabilité $1/2$. Seul compte le résultat du groupe, et un plan modifie la manière dont les bonnes et les mauvaises annonces se recoupent, pas leur nombre total.}
\iqfollow{Et avec sept personnes~?}{Une probabilité de victoire de $7/8$, en utilisant le code de Hamming de longueur 7~: il faut n'annoncer une couleur que si les chapeaux vus, avec une couleur donnée pour le sien, formeraient un mot du code, et annoncer alors l'autre couleur.}
\end{solution}

\begin{solution}{iq:iv1:logic-riddles:32}
Les 100 personnes aux yeux bleus partent ensemble la 100e nuit, et personne d'autre ne part. Le visiteur a rendu public un fait qui ne l'était pas auparavant.

Raisonnons par récurrence sur le nombre $k$ de personnes aux yeux bleus\index{connaissance commune}. Si $k = 1$, cette personne ne voit aucun œil bleu et part la nuit 1. Si $k = 2$, chacune voit une personne aux yeux bleus. Chacune pense~: «~Si je n'ai pas les yeux bleus, l'autre partira ce soir.~» L'autre ne part pas, donc les deux partent la nuit 2. En général, chacune des $k$ personnes aux yeux bleus en voit $k - 1$, et attend de voir si celles-ci partent la nuit $k - 1$. Quand ce n'est pas le cas, elles partent toutes la nuit $k$. Une personne aux yeux marron voit 100 yeux bleus, et lorsque ceux-ci partent la nuit 100, elle apprend seulement que ses propres yeux ne sont pas bleus, ce qui ne lui donne pas sa couleur.

Qu'y avait-il de nouveau~? Tout le monde savait que quelqu'un avait les yeux bleus. Mais avec deux personnes aux yeux bleus, A ne sait pas que B sait. Avec trois, A ne sait pas que B sait que C sait. Avec 100, la chaîne des «~sait que~» doit compter 100 niveaux, et elle s'arrête court sans une déclaration publique. L'annonce rend vrai à la fois chaque niveau. Une simulation par mondes possibles avec jusqu'à 5 personnes aux yeux bleus parmi 6 reproduit les jours.

\iqwrong{Dire que rien ne se passe, puisque le visiteur ne leur a rien appris qu'ils ne savaient déjà. Le premier niveau de connaissance n'a pas changé. Les niveaux supérieurs, ce que chacun sait de ce que les autres savent, ont changé.}
\iqwhy{Une information publique peut faire bouger un marché même si chacun la connaissait en privé. Si l'on ignore que les autres savent, on n'agit pas comme lorsqu'on le sait.}
\end{solution}

\begin{solution}{iq:iv1:logic-riddles:33}
L'un des prisonniers est le compteur. Chacun des autres allume la lumière la première fois qu'il la trouve éteinte, et jamais plus. Seul le compteur l'éteint, et ajoute 1 à chaque fois. Lorsque le compteur atteint 99, il annonce la fin. Cela prend environ $10\,418$ jours en moyenne, soit environ $28.5$ ans.

Cela ne peut pas échouer, car le compte n'augmente que lorsqu'un nouveau prisonnier a signalé son passage\index{protocole de comptage}. L'interrupteur est éteint au départ, donc il n'y a aucun faux compte depuis le début. Chacun des 99 non-compteurs allume la lumière au plus une fois.

Le temps est une somme de deux attentes pour chacun des 99 signaux. Si $100 - k$ non-compteurs n'ont pas encore signalé leur passage, l'un d'eux arrive en moyenne au bout de $\frac{100}{100 - k}$ jours. Il faut ensuite que le compteur passe, ce qui prend 100 jours en moyenne. D'où
\[
\sum_{k = 1}^{99} \Bigl(\frac{100}{100 - k} + 100\Bigr) = 100 H_{99} + 9900 \approx 10\,418.
\]
Presque tout ce temps est une attente du compteur. Une simulation du plan donne une moyenne à moins de $1.5~\%$ de ce résultat.

\iqwrong{Déclarer la fin après un grand nombre de jours, au motif que tout le monde doit être passé d'ici là. Cela peut devenir très probable, mais la question exige la certitude. Un décompte est sûr, une limite de temps ne l'est pas.}
\iqfollow{Et si vous ne connaissez pas l'état initial de l'interrupteur~?}{Alors il faut compter deux fois plus~: chaque non-compteur allume la lumière deux fois, et le compteur s'arrête à $2 \cdot 99 = 198$, ce qui prend à peu près deux fois plus de temps. Le faux compte initial ne peut pas à lui seul atteindre le total.}
\end{solution}

\begin{solution}{iq:iv1:logic-riddles:34}
Numérotons les cases de 0 à 63. Soit $P$ le ou exclusif des numéros de toutes les cases qui montrent face. Le prisonnier A retourne la pièce sur la case $P \oplus t$, où $t$ est la cible. Le prisonnier B calcule le nouveau $P$, qui vaut $t$. La méthode fonctionne exactement lorsque le nombre de cases est une puissance de 2.

Retourner la pièce sur la case $s$ remplace $P$ par $P \oplus s$\index{codage de parité}\index{ou exclusif}, que la pièce soit initialement face ou pile, car le ou exclusif d'un nombre avec lui-même vaut 0. A veut donc $P \oplus s = t$, c'est-à-dire $s = P \oplus t$. C'est toujours une case valide, car $P$ et $t$ ont tous deux 6 bits. Chaque cible est donc accessible à A, en exactement un retournement. Un test sur des plateaux et des cibles aléatoires ne révèle aucune erreur.

Ce qui est particulier avec 64~: le plan attribue une case à chaque disposition, et depuis n'importe quelle disposition, un seul retournement doit atteindre chaque cible. Sur le cube des dispositions à $n$ coordonnées, chaque sommet a $n$ voisins, et les $n$ couleurs doivent toutes y apparaître. Chaque couleur se trouve donc sur exactement $2^n/n$ sommets, et $n$ doit diviser $2^n$. Donc $n$ doit être une puissance de 2. Une recherche exhaustive pour trois cases ne trouve aucun coloriage, comme le prédit cet argument.

\iqwrong{Essayer une parité séparée sur les lignes et sur les colonnes, ce qui marche pour quelques bits mais ne peut pas en traiter six à la fois. Chaque bit de l'adresse doit être contrôlé par un ensemble de cases qui se recoupe avec les autres de la bonne façon, et le ou exclusif le permet.}
\iqwhy{C'est un code parfait. Un bit modifié porte exactement $\log_2 n$ bits d'information, et la même idée sous-tend le code de Hamming utilisé pour corriger les erreurs en mémoire.}
\end{solution}

\begin{solution}{iq:iv1:logic-riddles:35}
Parmi 5 cartes, deux ont la même couleur. L'assistant en cache une de ces deux et place l'autre en premier, pour fixer la couleur et un rang de base. Les trois autres cartes, dans l'un des 6 ordres possibles, indiquent la distance entre la base et la carte cachée. Un argument de comptage permet jusqu'à 124 cartes.

Il y a 13 rangs sur un cercle, donc l'une des deux cartes de même couleur est à au plus 6 pas dans le sens horaire de l'autre\index{principe des tiroirs}. Il faut montrer la carte à partir de laquelle on compte, et cacher l'autre. Ordonner les trois restantes, selon un ordre fixe de toutes les cartes, pour coder la distance $d$ de 1 à 6, car $3! = 6$. Le magicien lit la première carte, ajoute $d$ sur le cercle, et garde la couleur. Un test sur 40\,000 mains aléatoires retrouve la carte cachée à chaque fois.

Pour la borne de comptage avec un jeu de $n$ cartes, chaque main doit être envoyée vers une rangée différente de 4 cartes. Il y a $\binom{n}{5}$ mains, et $n(n - 1)(n - 2)(n - 3)$ rangées. Cela exige $\frac{n - 4}{120} \le 1$, donc $n \le 124$. Cette borne peut être atteinte~: chaque main peut utiliser 120 rangées (5 choix pour la carte cachée fois 4~! ordres), chaque rangée convient à $n - 4$ mains, et comme $n - 4 \le 120$, un couplage existe par le théorème de Hall\index{théorème de Hall du mariage}. Cette preuve ne donne pas de formule simple.

\iqwrong{Dire que 4~! = 24 ordres ne peuvent pas coder une carte parmi 48. L'assistant choisit aussi la carte à cacher, ce qui donne 5 fois plus de liberté. Avec la couleur donnée par la première carte, il ne faut qu'une petite partie de cette liberté.}
\end{solution}
